\scanpage{8}\mspage{7}
\[
\begin{aligned}
\leq,\quad\ll,\quad\preccurlyeq&\qquad\text{order relations}\\
&\qquad\text{\unsure{in|of} }\unsure{\mathcal{F}|\mathfrak{F}}
\end{aligned}\tag*{(1.10).}
\]
and the relations induced in $\mathcal{M}$
\[
\leq,\quad\ll\quad\text{in }\mathcal{M}\tag*{(1.11)}
\]
% slanted
\margin{\uline{NB} \unsure{The|the} order relations induced by $\leq$, $\ll$ on $\mathcal{L}$ are discrete.}
[nb the order relation $\preccurlyeq$ induced in $\mathcal{M}$ is less important, as a specific structure in $\mathcal{M}$].

% slanted
\margin{\uline{NB} $G\leq F\Longrightarrow G\ll F$ : a sub-figure of $F$ refines $F$}
\begin{center}
\begin{tabular}{@{}l@{\quad}p{.64\linewidth}@{}}
$\deleted{F}G\deleted{\unsure{\ll|<}}\leq\deleted{G}F$ & $\deleted{F}G$ is called a \uline{sub}-\uline{figure} of $\deleted{G}F$\\
& (If $\deleted{F}G\inserted{=X}\in\mathcal{M}$, it is therefore a \uline{stratum} of $\deleted{G}F$)\\[.45em]
$G\ll F$ & $G$ is called a \uline{refinement} of $F$\\
& (if $G=\unsure{x|X}\in\mathcal{M}$, it is therefore an element of the large \unsure{shadow|\deleted{o}shadow|\deleted{m}shadow} of $F$)\\[.45em]
$G\preccurlyeq F$ & $G$ is called a \uline{subdivision} of $F$,\\
& one says that the figure $G$ subdivides the figure $F$.
\end{tabular}
\end{center}

I note finally tw\deleted{\unsure{q|x}}o other important relations \deleted{in} between figures:
\[
\left\{\begin{aligned}
F\text{ and }G&\quad\text{\uline{compatible}}\\
F\text{ and }G&\quad\text{\uline{disjoint}}\qquad,
\end{aligned}\right.\tag*{(1.12)}
\]
both of which can be expressed in terms of the order relation $\leq$:

$F$ and $G$ compatible: they are subfigures of a $\overline m$ figure, \unsure{i..e|i.e} $\{F,G\}$ bounded above for $\leq$;

$F$ and $G$ disjoint: $F$ and $G$ compatible, and their
